% =====================================================================
%  COURS DE CHIMIE — FICHE 7
%  Loi des gaz
%  Document généré en LaTeX avec figures TikZ tracées « from scratch ».
%  Compilation : pdflatex (2 passes pour la table des matières)
%               ou  latexmk -pdf Cours_Fiche-07_Loi-des-gaz.tex
% =====================================================================
\documentclass[11pt,a4paper]{article}

% ---------- Encodage & langue ----------
\usepackage[utf8]{inputenc}
\usepackage[T1]{fontenc}
\usepackage{lmodern}
\usepackage[french]{babel}

% ---------- Mathématiques ----------
\usepackage{amsmath,amssymb,mathtools}
\usepackage{icomma}              % virgule décimale correcte en mode maths

% ---------- Unités & chimie ----------
\usepackage{siunitx}
\sisetup{output-decimal-marker={,},per-mode=symbol}
% unités non SI utiles pour les gaz
\DeclareSIUnit{\atm}{atm}
\DeclareSIUnit{\bar}{bar}
\DeclareSIUnit{\mmHg}{mmHg}
\DeclareSIUnit{\torr}{Torr}
\usepackage[version=4]{mhchem}   % formules chimiques : \ce{N2}, \ce{CO2}...

% ---------- Couleurs, boîtes, graphismes ----------
\usepackage{xcolor}
\usepackage{colortbl}   % \rowcolor dans les tableaux
\usepackage{tikz}
\usepackage{pgfplots}
\usepackage[most]{tcolorbox}
\usepackage{enumitem}
\usepackage{array}
\usepackage{booktabs}
\usepackage{multicol}
\usepackage{fancyhdr}
\usepackage[hidelinks]{hyperref}

\usetikzlibrary{arrows.meta,positioning,calc,decorations.pathmorphing,
  decorations.markings,angles,quotes,patterns,backgrounds,
  shapes.geometric,shapes.misc,fadings,intersections,fit}
\pgfplotsset{compat=1.18}

% ---------- Géométrie de page ----------
\usepackage[a4paper,margin=2.2cm,headheight=26pt,headsep=10pt]{geometry}

% =====================================================================
%  PALETTE DE COULEURS  (mauve = couleur de la section Chimie du livre)
% =====================================================================
\definecolor{primary}{RGB}{150,98,162}      % mauve (couleur du livre — Chimie)
\definecolor{primarydark}{RGB}{92,52,110}
\definecolor{defcol}{RGB}{0,114,178}        % bleu  — définitions
\definecolor{thmcol}{RGB}{0,140,120}        % vert-bleu — théorèmes
\definecolor{formulacol}{RGB}{198,93,9}     % orange — formules
\definecolor{reglecol}{RGB}{150,98,162}     % mauve — règles
\definecolor{remarkcol}{RGB}{120,94,200}    % violet — remarques
\definecolor{attncol}{RGB}{200,40,55}       % rouge — pièges
\definecolor{excol}{RGB}{0,135,90}          % vert — exemples
\definecolor{methcol}{RGB}{90,90,100}       % gris — méthode
% --- couleurs spécifiques « gaz » ---
\definecolor{gazcol}{RGB}{0,120,190}        % bleu  — molécules de gaz
\definecolor{pcol}{RGB}{200,40,55}          % rouge — pression
\definecolor{vcol}{RGB}{0,135,90}           % vert  — volume
\definecolor{tcol}{RGB}{198,93,9}           % orange — température
\definecolor{ncol}{RGB}{120,94,200}         % violet — quantité de matière
\definecolor{azote}{RGB}{60,90,200}         % bleu  — diazote
\definecolor{oxygene}{RGB}{210,60,60}       % rouge — dioxygène

% =====================================================================
%  STYLE COMMUN DES BOÎTES
% =====================================================================
\tcbset{
  boxbase/.style={enhanced,breakable,boxrule=0.9pt,arc=2.5pt,
     left=3mm,right=3mm,top=2.3mm,bottom=2.3mm,
     fonttitle=\bfseries,coltitle=white,
     toptitle=0.6mm,bottomtitle=0.6mm},
}

% ---- Définition (numérotée) ----
\newtcolorbox[auto counter,number within=section]{definition}[1][]{%
  boxbase,colback=defcol!5,colframe=defcol,
  title={\textsc{Définition}~\thetcbcounter\ \ifx\relax#1\relax\relax\else\textmd{--- \itshape #1}\fi}}

% ---- Théorème / Propriété (numéroté) ----
\newtcolorbox[auto counter,number within=section]{theoreme}[1][]{%
  boxbase,colback=thmcol!6,colframe=thmcol,
  title={\textsc{Loi / Propriété}~\thetcbcounter\ \ifx\relax#1\relax\relax\else\textmd{--- \itshape #1}\fi}}

% ---- Formule (numérotée) ----
\newtcolorbox[auto counter,number within=section]{formule}[1][]{%
  boxbase,colback=formulacol!6,colframe=formulacol,
  title={\textsc{Formule}~\thetcbcounter\ \ifx\relax#1\relax\relax\else\textmd{--- \itshape #1}\fi}}

% ---- Remarque ----
\newtcolorbox{remarque}[1][]{%
  boxbase,colback=remarkcol!6,colframe=remarkcol,
  title={\textsc{Remarque}\ifx\relax#1\relax\relax\else\textmd{ : \itshape #1}\fi}}

% ---- Attention / piège ----
\newtcolorbox{attention}[1][]{%
  boxbase,colback=attncol!6,colframe=attncol,
  title={\textsc{Attention}\ifx\relax#1\relax\relax\else\textmd{ : \itshape #1}\fi}}

% ---- Exemple ----
\newtcolorbox{exemple}[1][]{%
  boxbase,colback=excol!5,colframe=excol,
  title={\textsc{Exemple}\ifx\relax#1\relax\relax\else\textmd{ : \itshape #1}\fi}}

% ---- Méthode ----
\newtcolorbox{methode}[1][]{%
  boxbase,colback=methcol!7,colframe=methcol,sharp corners,
  title={\textsc{Méthode}\ifx\relax#1\relax\relax\else\textmd{ : \itshape #1}\fi}}

% =====================================================================
%  ENVIRONNEMENT « où : »  (légende des grandeurs d'une formule)
% =====================================================================
\newenvironment{ou}{%
  \par\smallskip\noindent\textbf{où :}\nopagebreak
  \begin{itemize}[leftmargin=1.8em,topsep=2pt,itemsep=1.5pt,parsep=0pt]}%
  {\end{itemize}}

% =====================================================================
%  EN-TÊTES ET PIEDS DE PAGE
% =====================================================================
\pagestyle{fancy}
\fancyhf{}
\fancyhead[L]{\small\textcolor{primarydark}{\textbf{Fiche 7} — Loi des gaz}}
\fancyhead[R]{\small\textcolor{primarydark}{\textsc{Chimie}}}
\fancyfoot[C]{\small\thepage}
\renewcommand{\headrulewidth}{0.8pt}
\renewcommand{\headrule}{\hbox to\headwidth{\color{primary}\leaders\hrule height \headrulewidth\hfill}}

% Titres de section en couleur
\usepackage{titlesec}
\titleformat{\section}{\Large\bfseries\color{primarydark}}{\thesection}{0.6em}{}
\titleformat{\subsection}{\large\bfseries\color{primary}}{\thesubsection}{0.6em}{}
\titleformat{\subsubsection}{\normalsize\bfseries\color{primary!80!black}}{\thesubsubsection}{0.6em}{}

% Raccourcis
\newcommand{\gp}{gaz parfait}
\newcommand{\GP}{gaz parfaits}
\newcommand{\Vm}{V_{\mathrm{m}}}

% --- macro pour dessiner un « nuage » de molécules de gaz dans un rectangle ---
% #1 = nombre de points, #2..#5 = xmin,ymin,xmax,ymax (utilisé manuellement)

% =====================================================================
\begin{document}
\sloppy

% ---------------------------------------------------------------------
%  PAGE DE TITRE
% ---------------------------------------------------------------------
\begingroup
\thispagestyle{empty}
\centering
\vspace*{1.2cm}

\begin{tikzpicture}
\node[rectangle,rounded corners=8pt,fill=primary,text=white,
      minimum width=15.5cm,minimum height=2.6cm,align=center,inner sep=10pt]
  {{\Huge\bfseries Loi des gaz}\\[6pt]
   {\large\bfseries Du gaz parfait aux pressions partielles}};
\end{tikzpicture}

\vspace{0.3cm}
{\large\color{primarydark}\textsc{Cours de chimie — Fiche n\textsuperscript{o}\,7}}

\vspace{0.7cm}

% --- Illustration de couverture : piston + molécules + équation d'état ---
\begin{tikzpicture}[scale=1.0]
  % cylindre
  \draw[thick] (0,0) rectangle (4.2,3);
  \fill[gazcol!8] (0,0) rectangle (4.2,2.3);
  % piston
  \draw[fill=primary!30,thick] (0,2.3) rectangle (4.2,2.65);
  \draw[thick] (2.0,2.65) -- (2.0,3.5);
  \draw[thick,fill=primary!30] (1.5,3.5) rectangle (2.7,3.75);
  \node[above] at (2.1,3.75) {\footnotesize piston mobile};
  % molécules (points) avec petites flèches de vitesse
  \foreach \x/\y in {0.6/0.5,1.4/1.1,2.3/0.4,3.1/1.4,3.7/0.7,0.9/1.7,
                     2.0/1.9,2.8/0.9,1.2/0.3,3.4/1.9,0.4/1.2,3.9/1.6}{
     \fill[gazcol] (\x,\y) circle (2.4pt);
     \draw[gazcol,->,>=Stealth,thin] (\x,\y) -- ++(0.32,0.18);}
  % flèche pression sur les parois
  \foreach \y in {0.5,1.2,1.9}{
     \draw[pcol,->,>=Stealth] (0.05,\y) -- (-0.45,\y);
     \draw[pcol,->,>=Stealth] (4.15,\y) -- (4.65,\y);}
  \node[pcol,font=\scriptsize,rotate=90] at (-0.75,1.2) {pression $p$};
  % equation
  \node[font=\Large] at (7.6,1.7)
     {$\textcolor{pcol}{p}\,\textcolor{vcol}{V}=\textcolor{ncol}{n}\,R\,\textcolor{tcol}{T}$};
  \node[font=\scriptsize,align=left,anchor=west] at (6.05,0.7)
     {\textcolor{pcol}{$p$} : pression\\
      \textcolor{vcol}{$V$} : volume\\
      \textcolor{ncol}{$n$} : quantité de matière\\
      \textcolor{tcol}{$T$} : température};
\end{tikzpicture}

\vspace{0.2cm}

\begin{tcolorbox}[boxbase,colback=primary!5,colframe=primary,width=15cm,
    title={\textsc{Objectifs pédagogiques de la fiche}}]
À l'issue de ce cours, l'étudiant·e sera capable de :
\begin{itemize}[leftmargin=1.6em,topsep=2pt,itemsep=1pt]
  \item décrire l'\textbf{état gazeux} à l'aide des quatre variables d'état $p$, $V$, $n$, $T$ et maîtriser leurs \textbf{unités} ;
  \item énoncer et appliquer les trois \textbf{lois historiques} : Boyle-Mariotte, Charles--Gay-Lussac et Avogadro ;
  \item utiliser l'\textbf{équation d'état des \GP} $pV=nRT$ et la valeur de la constante $R$ ;
  \item calculer la \textbf{masse volumique} et la \textbf{masse molaire} d'un gaz, ainsi que sa \textbf{densité} par rapport à l'air ;
  \item exploiter le \textbf{volume molaire} $\qty{22{,}4}{\litre\per\mole}$ aux conditions normales (CNTP) ;
  \item traiter un \textbf{mélange gazeux} avec la \textbf{loi de Dalton} (pressions partielles, fraction molaire).
\end{itemize}
\end{tcolorbox}

\vspace{0.4cm}
{\small\color{black!60} Document de synthèse rédigé d'après la Fiche 7 du livre — figures originales tracées en TikZ.}
\par
\endgroup

% ---------------------------------------------------------------------
%  TABLE DES MATIÈRES
% ---------------------------------------------------------------------
\newpage
\renewcommand{\contentsname}{\textcolor{primarydark}{Table des matières}}
\tableofcontents
\thispagestyle{fancy}
\newpage

% =====================================================================
\section{L'état gazeux : décrire un gaz par quatre nombres}
% =====================================================================
La chimie étudie souvent des réactions qui produisent ou consomment des \textbf{gaz} (combustions,
fermentations, synthèses industrielles comme celle de l'ammoniac \ce{NH3}). Contrairement à un solide
ou à un liquide, un gaz n'a ni forme ni volume propres : il \emph{occupe tout le récipient} qui le
contient. Pour le décrire complètement, il suffit pourtant de \textbf{quatre grandeurs}, appelées
\emph{variables d'état}.

\begin{definition}[gaz et variables d'état]
Un \textbf{gaz} est un état de la matière dans lequel les molécules sont très espacées, animées d'une
agitation incessante et désordonnée, et n'interagissent presque pas entre elles. L'état d'un
échantillon de gaz est entièrement fixé par la donnée de \textbf{quatre variables d'état} :
\begin{itemize}[leftmargin=1.6em,topsep=2pt,itemsep=1.5pt]
  \item la \textcolor{pcol}{\textbf{pression}} $p$ ;
  \item le \textcolor{vcol}{\textbf{volume}} $V$ ;
  \item la \textcolor{ncol}{\textbf{quantité de matière}} $n$ (nombre de moles) ;
  \item la \textcolor{tcol}{\textbf{température}} $T$.
\end{itemize}
Les lois des gaz sont précisément les \emph{relations} qui lient ces quatre nombres entre eux.
\end{definition}

\subsection{Le modèle microscopique : d'où vient la pression ?}

À l'échelle microscopique, un gaz est un essaim de molécules qui se déplacent en ligne droite et
rebondissent sur les parois du récipient. Chaque choc d'une molécule contre la paroi exerce une
minuscule poussée ; la somme de ces milliards de chocs par seconde et par unité de surface produit
une force continue : c'est la \textbf{pression}.

\begin{center}
\begin{tikzpicture}[scale=1.0]
  % boîte
  \draw[thick] (0,0) rectangle (5,3.2);
  \fill[gazcol!6] (0,0) rectangle (5,3.2);
  % molécules + vecteurs vitesse
  \foreach \x/\y/\dx/\dy in {0.8/0.7/0.4/0.3, 1.7/2.3/-0.3/0.35, 2.6/1.2/0.45/-0.2,
        3.5/2.6/0.2/0.4, 4.2/0.9/-0.4/0.25, 1.2/1.8/0.35/-0.35, 3.0/0.5/-0.25/0.4,
        4.4/2.2/0.3/0.3, 2.1/2.8/0.4/-0.2, 3.7/1.6/-0.4/-0.3}{
     \fill[gazcol] (\x,\y) circle (2.6pt);
     \draw[gazcol,->,>=Stealth,thin] (\x,\y) -- ++(\dx,\dy);}
  % choc sur la paroi droite illustré
  \fill[pcol] (4.85,1.5) circle (2.6pt);
  \draw[pcol,->,>=Stealth,thick] (4.85,1.5) -- (5.35,1.5);
  \node[pcol,right,font=\scriptsize] at (5.4,1.5) {force du choc sur la paroi};
  % légende
  \node[gazcol,font=\scriptsize] at (2.5,-0.4) {molécules en mouvement désordonné (agitation thermique)};
\end{tikzpicture}
\end{center}

\noindent Ce modèle explique immédiatement deux faits :
\begin{itemize}[leftmargin=1.6em,topsep=2pt,itemsep=1.5pt]
  \item plus on \textbf{chauffe} le gaz, plus les molécules vont vite, plus les chocs sont violents et
        fréquents $\Rightarrow$ la pression (ou le volume) \emph{augmente} ;
  \item plus on \textbf{comprime} le gaz (volume plus petit), plus les molécules frappent souvent les
        parois $\Rightarrow$ la pression \emph{augmente}.
\end{itemize}

\subsection{Le gaz parfait : un modèle idéal}

\begin{definition}[gaz parfait]
Un \textbf{\gp} est un gaz idéalisé qui obéit exactement aux hypothèses suivantes :
\begin{enumerate}[leftmargin=1.8em,topsep=2pt,itemsep=1.5pt,label=(\arabic*)]
  \item les molécules sont \textbf{ponctuelles} : leur volume propre est négligeable devant le volume
        offert ;
  \item elles \textbf{n'interagissent pas} à distance (pas de forces d'attraction ni de répulsion entre
        elles) ;
  \item leurs \textbf{chocs} (entre elles et sur les parois) sont \emph{élastiques} : ils ne dissipent
        pas d'énergie.
\end{enumerate}
\end{definition}

\begin{remarque}
Le modèle du \gp\ est une \emph{idéalisation}. Les gaz réels s'en rapprochent d'autant mieux qu'ils
sont \textbf{peu denses} : basse pression et température pas trop basse. Dans les conditions usuelles
du laboratoire, l'air, le diazote, le dioxygène, etc.\ se comportent en très bonne approximation comme
des \GP. Sauf mention contraire, \textbf{tous les gaz de ce cours sont supposés parfaits}.
\end{remarque}

\subsection{Les quatre grandeurs et leurs unités}

Maîtriser les \textbf{unités} est ici décisif : la valeur de la constante $R$ (vue plus loin) dépend du
choix d'unités, et une erreur d'unité est la première source de fautes dans les calculs de gaz.

% ----- Pression -----
\subsubsection{La pression \texorpdfstring{$p$}{p}}

\begin{definition}[pression]
La \textbf{pression} est la force pressante exercée perpendiculairement à une surface, rapportée à
l'aire de cette surface :
\[ p=\frac{F}{S}. \]
\end{definition}

\begin{ou}
  \item $p$ : la pression, en pascals (\si{\pascal}) dans le Système international ;
  \item $F$ : la force pressante, en newtons (\si{\newton}) ;
  \item $S$ : l'aire de la surface pressée, en mètres carrés (\si{\meter\squared}).
\end{ou}

\noindent L'unité SI est le \textbf{pascal} : $\qty{1}{\pascal}=\qty{1}{\newton\per\meter\squared}$.
C'est une très petite unité, on utilise donc souvent ses multiples et des unités usuelles. Le tableau
suivant rassemble les conversions à connaître.

\begin{center}\renewcommand{\arraystretch}{1.35}
\begin{tabular}{l l l}
\toprule
\rowcolor{primary!12}
\textbf{Unité} & \textbf{Symbole} & \textbf{Équivalence}\\
\midrule
pascal              & \si{\pascal}   & unité SI de base\\
kilopascal          & \si{\kilo\pascal} & $\qty{1}{\kilo\pascal}=\qty{1000}{\pascal}$\\
bar                 & \si{\bar}      & $\qty{1}{\bar}=\qty{e5}{\pascal}=\qty{100}{\kilo\pascal}$\\
atmosphère          & \si{\atm}      & $\qty{1}{\atm}=\qty{101325}{\pascal}\approx\qty{101{,}3}{\kilo\pascal}$\\
millimètre de mercure & \si{\mmHg}   & $\qty{1}{\atm}=\qty{760}{\mmHg}$\\
\bottomrule
\end{tabular}
\end{center}

\begin{attention}
$\qty{1}{\atm}\approx\qty{1}{\bar}$ mais \emph{ce n'est pas exactement égal} :
$\qty{1}{\atm}=\qty{1{,}013}{\bar}$. Dans les calculs de la fiche, on emploie surtout l'atmosphère
($\qty{1}{\atm}$) car la valeur de $R$ y devient très simple.
\end{attention}

% ----- Volume -----
\subsubsection{Le volume \texorpdfstring{$V$}{V}}

\begin{definition}[volume]
Le \textbf{volume} $V$ est l'espace occupé par le gaz, c'est-à-dire le volume du récipient (le gaz
remplit tout l'espace disponible).
\end{definition}

\begin{ou}
  \item $V$ : le volume, en mètres cubes (\si{\meter\cubed}) dans le SI, mais le plus souvent en
        litres (\si{\litre}) en chimie.
\end{ou}

\noindent Conversions indispensables :
\[ \qty{1}{\litre}=\qty{1}{\deci\meter\cubed}=\qty{e-3}{\meter\cubed}=\qty{1000}{\milli\litre},
   \qquad \qty{1}{\meter\cubed}=\qty{1000}{\litre}. \]
On retiendra notamment $\qty{1}{\milli\litre}=\qty{1}{\centi\meter\cubed}$ et
$\qty{1}{\meter\cubed}=\qty{1000}{\litre}$, très utile pour passer du volume d'une pièce (en \si{\meter\cubed})
à un volume de gaz (en \si{\litre}).

% ----- Température -----
\subsubsection{La température \texorpdfstring{$T$}{T}}

\begin{definition}[température absolue]
Dans les lois des gaz, la température doit \emph{toujours} être exprimée en \textbf{kelvins}
(\si{\kelvin}), unité de la \textbf{température absolue}. Le zéro de cette échelle, $\qty{0}{\kelvin}$,
est le \textbf{zéro absolu} : la température la plus basse possible, où l'agitation des molécules cesse.
\end{definition}

\begin{formule}[conversion Celsius $\leftrightarrow$ Kelvin]
\[ \boxed{\,T(\si{\kelvin}) = \theta(\si{\degreeCelsius}) + 273{,}15\,} \]
\end{formule}

\begin{ou}
  \item $T$ : la température absolue, en kelvins (\si{\kelvin}) ;
  \item $\theta$ : la température en degrés Celsius (\si{\degreeCelsius}).
\end{ou}

\noindent En pratique, dans cette fiche, on arrondit fréquemment à $T=\theta+273$. Une variation de
$\qty{1}{\degreeCelsius}$ correspond exactement à une variation de $\qty{1}{\kelvin}$ : seules les
\emph{origines} des deux échelles diffèrent.

\begin{center}
\begin{tikzpicture}[scale=1.0]
  % axe horizontal Celsius
  \draw[thick,->,>=Stealth] (-0.5,0) -- (10.5,0);
  \draw[thick,->,>=Stealth] (-0.5,1.4) -- (10.5,1.4);
  % graduations communes : -273.15, 0, 100, 273
  % échelle : x = 1 + (theta)/40 ... on place manuellement
  % positions : -273 -> 0 ; 0 -> 6.83 ; 27 -> 7.5 ; 100 -> 9.3  (1 unité x = 40°C)
  \foreach \xc/\cel/\kel in {0/-273/0, 6.83/0/273, 7.5/27/300, 9.27/100/373}{
     \draw (\xc,1.25) -- (\xc,1.55);
     \draw (\xc,-0.15) -- (\xc,0.15);
     \node[tcol,font=\scriptsize] at (\xc,1.85) {\kel\,K};
     \node[font=\scriptsize] at (\xc,-0.45) {\cel\,\textdegree C};}
  \node[left,tcol,font=\small\bfseries] at (-0.5,1.4) {K};
  \node[left,font=\small\bfseries] at (-0.5,0) {\textdegree C};
  % zéro absolu
  \draw[densely dashed,attncol] (0,-0.6) -- (0,2.1);
  \node[attncol,font=\scriptsize,align=center,above] at (0,2.1) {zéro absolu};
  % décalage 273.15
  \draw[<->,>=Stealth,primary] (0,0.7) -- (6.83,0.7);
  \node[primary,font=\scriptsize,fill=white] at (3.4,0.7) {décalage de $273{,}15$};
\end{tikzpicture}
\end{center}

\begin{attention}
Oublier de convertir en kelvins est l'erreur \textbf{numéro un}. Une réaction « à $\qty{27}{\degreeCelsius}$ »
se calcule avec $T=\qty{300}{\kelvin}$, \emph{jamais} avec $27$. Un rapport $T_2/T_1$ calculé en degrés
Celsius donne un résultat totalement faux (on ne peut pas diviser des températures Celsius).
\end{attention}

% ----- Quantité de matière -----
\subsubsection{La quantité de matière \texorpdfstring{$n$}{n}}

\begin{definition}[mole et quantité de matière]
La \textbf{quantité de matière} $n$ mesure le \emph{nombre de paquets} de molécules, exprimé en
\textbf{moles} (\si{\mole}). Une mole contient $N_A=\qty{6{,}022e23}{}$ entités (nombre d'Avogadro). On
la relie à la masse $m$ et à la masse molaire $M$ par :
\[ n=\frac{m}{M}. \]
\end{definition}

\begin{ou}
  \item $n$ : la quantité de matière, en moles (\si{\mole}) ;
  \item $m$ : la masse de gaz, en grammes (\si{\gram}) ;
  \item $M$ : la masse molaire du gaz, en grammes par mole (\si{\gram\per\mole}).
\end{ou}

\begin{remarque}
Cette relation $n=m/M$ a été étudiée à la Fiche 4. Elle est essentielle ici : c'est le pont entre la
\emph{masse} d'un gaz (ce que l'on pèse) et sa \emph{quantité de matière} (ce qui intervient dans
$pV=nRT$).
\end{remarque}

% =====================================================================
\section{Les trois lois historiques des gaz}
% =====================================================================
Avant de disposer de l'équation générale $pV=nRT$, les chimistes et physiciens des
\textsc{xvii\textsuperscript{e}}--\textsc{xix\textsuperscript{e}} siècles ont découvert
\emph{expérimentalement} trois lois partielles. Chacune fixe \textbf{deux} des quatre variables et
décrit comment les \textbf{deux} autres sont liées. Ce sont les briques de base de la loi des \GP.

\begin{center}\renewcommand{\arraystretch}{1.4}
\begin{tabular}{l l l l}
\toprule
\rowcolor{primary!12}
\textbf{Loi} & \textbf{Ce qui est fixé} & \textbf{Ce qui varie} & \textbf{Relation}\\
\midrule
Boyle-Mariotte       & $n,\ T$ & $p,\ V$ & $pV=\text{cste}$\\
Charles--Gay-Lussac  & $n,\ p$ & $V,\ T$ & $V/T=\text{cste}$\\
Avogadro             & $p,\ T$ & $V,\ n$ & $V/n=\text{cste}$\\
\bottomrule
\end{tabular}
\end{center}

% =====================================================================
\subsection{Loi de Boyle-Mariotte (transformation à \texorpdfstring{$T$}{T} constante)}

\begin{theoreme}[loi de Boyle-Mariotte]
Pour une quantité fixée de gaz ($n$ constant) maintenue à \textbf{température constante} ($T$
constante), le produit de la pression par le volume reste \textbf{constant}. Autrement dit, $p$ et $V$
sont \textbf{inversement proportionnels} :
\[ \boxed{\,p_1\times V_1 = p_2\times V_2\,} \]
Si l'on comprime le gaz (on diminue $V$), sa pression $p$ augmente dans la même proportion, et
réciproquement.
\end{theoreme}

\begin{ou}
  \item $p_1,\ V_1$ : pression et volume dans l'état initial ;
  \item $p_2,\ V_2$ : pression et volume dans l'état final ;
  \item les unités de $p$ et de $V$ sont libres, à condition d'être \emph{les mêmes} de chaque côté de
        l'égalité (par ex.\ $p$ en \si{\atm} et $V$ en \si{\litre}).
\end{ou}

\noindent Une transformation à température constante s'appelle une transformation \textbf{isotherme}.
Dans un diagramme $(V,p)$, l'ensemble des états accessibles est une \textbf{hyperbole}
$p=\text{cste}/V$.

\begin{center}
\begin{tikzpicture}
\begin{axis}[
    width=9.5cm,height=6cm,
    xlabel={Volume $V$ (\si{\litre})},
    ylabel={Pression $p$ (\si{\atm})},
    xmin=0,xmax=10,ymin=0,ymax=7,
    axis lines=left,
    xtick={2,4,6,8},ytick={1,2,3,4,5,6},
    every axis x label/.style={at={(current axis.right of origin)},anchor=west},
    every axis y label/.style={at={(current axis.above origin)},anchor=south},
    clip=false]
  % isotherme pV = 12
  \addplot[domain=2:9.5,samples=120,thick,gazcol]{12/x};
  \node[gazcol,anchor=west] at (axis cs:5.4,4.6) {isotherme $pV=12$};
  % deux états
  \addplot[only marks,mark=*,pcol,mark size=2.2pt] coordinates {(2,6) (6,2)};
  \node[pcol,anchor=west,font=\footnotesize] at (axis cs:2.2,6.2)
     {$(V_1,p_1)=(\qty{2}{\litre},\qty{6}{\atm})$};
  \node[pcol,anchor=west,font=\footnotesize] at (axis cs:6.2,2.3)
     {$(V_2,p_2)=(\qty{6}{\litre},\qty{2}{\atm})$};
  % pointillés
  \draw[densely dashed,gray] (axis cs:2,0)--(axis cs:2,6)--(axis cs:0,6);
  \draw[densely dashed,gray] (axis cs:6,0)--(axis cs:6,2)--(axis cs:0,2);
\end{axis}
\end{tikzpicture}
\end{center}

\noindent On vérifie sur le graphique : $p_1V_1=6\times2=\num{12}$ et $p_2V_2=2\times6=\num{12}$ : le
produit $pV$ est bien conservé.

\begin{exemple}[compression d'un gaz dans une seringue]
On enferme $\qty{24}{\litre}$ d'air dans un cylindre à la pression $p_1=\qty{1}{\atm}$, à température
ambiante. On pousse le piston jusqu'à porter la pression à $p_2=\qty{3}{\atm}$, \emph{sans changer la
température}. Quel est le nouveau volume $V_2$ ?

\smallskip
\textbf{Données et choix de la loi.} La température et la quantité de gaz ne changent pas : on est à
$n$ et $T$ constants. C'est exactement le cadre de Boyle-Mariotte :
\[ p_1 V_1 = p_2 V_2. \]

\textbf{Isolation de l'inconnue.} On cherche $V_2$, donc :
\[ V_2 = \frac{p_1 V_1}{p_2}. \]

\textbf{Application numérique.}
\[ V_2 = \frac{\qty{1}{\atm}\times\qty{24}{\litre}}{\qty{3}{\atm}}
       = \qty{8}{\litre}. \]

\textbf{Interprétation.} En triplant la pression, on a divisé le volume par $3$ ($24\to8$). C'est
cohérent avec la proportionnalité \emph{inverse} : $p$ et $V$ varient en sens opposés, et leur produit
$pV=\qty{24}{\atm\litre}$ est conservé.
\end{exemple}

% =====================================================================
\subsection{Loi de Charles--Gay-Lussac (transformation à \texorpdfstring{$p$}{p} constante)}

\begin{theoreme}[loi de Charles--Gay-Lussac]
Pour une quantité fixée de gaz ($n$ constant) maintenue à \textbf{pression constante} ($p$ constante),
le volume est \textbf{directement proportionnel à la température absolue} :
\[ \boxed{\,\frac{V_1}{T_1}=\frac{V_2}{T_2}\,} \]
Chauffer le gaz le dilate ; le refroidir le contracte, en proportion exacte de la température en
kelvins.
\end{theoreme}

\begin{ou}
  \item $V_1,\ T_1$ : volume et température absolue dans l'état initial ;
  \item $V_2,\ T_2$ : volume et température absolue dans l'état final ;
  \item $T_1,\ T_2$ \emph{obligatoirement} en kelvins (\si{\kelvin}) ; $V$ dans une unité quelconque,
        identique des deux côtés.
\end{ou}

\noindent Une transformation à pression constante s'appelle une transformation \textbf{isobare}. Dans
un diagramme $(T,V)$, les états forment une \textbf{droite passant par l'origine} $T=\qty{0}{\kelvin}$.
Le prolongement de cette droite vers les volumes nuls pointe précisément vers le \textbf{zéro absolu}
$\theta=\qty{-273{,}15}{\degreeCelsius}$ : c'est l'une des manières historiques de le mettre en
évidence.

\begin{center}
\begin{tikzpicture}
\begin{axis}[
    width=10.5cm,height=5.6cm,
    xlabel={Température $T$ (\si{\kelvin})},
    ylabel={Volume $V$ (\si{\litre})},
    xmin=-30,xmax=420,ymin=0,ymax=3.4,
    axis lines=left,
    xtick={0,100,200,273,300,400},
    xticklabels={$0$,$100$,$200$,$273$,$300$,$400$},
    ytick={1,2,3},
    every axis x label/.style={at={(current axis.right of origin)},anchor=west},
    every axis y label/.style={at={(current axis.above origin)},anchor=south},
    clip=false]
  % droite V = (V/T)*T avec V/T = 2/300
  \addplot[domain=0:400,samples=2,thick,vcol]{2/300*x};
  \node[vcol,anchor=west,font=\footnotesize] at (axis cs:250,2.55) {$V\propto T$ (isobare)};
  % états
  \addplot[only marks,mark=*,pcol,mark size=2.2pt] coordinates {(300,2) (360,2.4)};
  \draw[densely dashed,gray](axis cs:300,0)--(axis cs:300,2)--(axis cs:0,2);
  \draw[densely dashed,gray](axis cs:360,0)--(axis cs:360,2.4)--(axis cs:0,2.4);
  \node[pcol,anchor=west,font=\scriptsize] at (axis cs:305,1.75){$(\qty{300}{\kelvin},\qty{2}{\litre})$};
  \node[pcol,anchor=west,font=\scriptsize] at (axis cs:300,2.95){$(\qty{360}{\kelvin},\qty{2{,}4}{\litre})$};
  % zéro absolu
  \node[attncol,font=\scriptsize,anchor=south] at (axis cs:0,0.05){$0$ K = zéro absolu};
\end{axis}
\end{tikzpicture}
\end{center}

\begin{exemple}[un ballon que l'on chauffe]
Un ballon souple contient $V_1=\qty{2{,}0}{\litre}$ d'air à $\theta_1=\qty{27}{\degreeCelsius}$, à la
pression atmosphérique. On le place dans une étuve à $\theta_2=\qty{87}{\degreeCelsius}$ ; la pression
reste celle de l'atmosphère. Quel volume occupe-t-il alors ?

\smallskip
\textbf{Étape 1 — convertir en kelvins (indispensable).}
\[ T_1 = 27+273 = \qty{300}{\kelvin}, \qquad T_2 = 87+273 = \qty{360}{\kelvin}. \]

\textbf{Étape 2 — choisir la loi.} Pression constante, quantité de gaz constante $\Rightarrow$
Charles--Gay-Lussac :
\[ \frac{V_1}{T_1}=\frac{V_2}{T_2}
   \quad\Longrightarrow\quad
   V_2 = V_1\times\frac{T_2}{T_1}. \]

\textbf{Étape 3 — application numérique.}
\[ V_2 = \qty{2{,}0}{\litre}\times\frac{\qty{360}{\kelvin}}{\qty{300}{\kelvin}}
       = \qty{2{,}0}{\litre}\times1{,}2 = \qty{2{,}4}{\litre}. \]

\textbf{Interprétation.} La température absolue a augmenté de $20\,\%$ ($300\to360$), donc le volume
augmente exactement de $20\,\%$ ($2{,}0\to2{,}4$). \textbf{Piège évité :} si l'on avait gardé les
degrés Celsius, on aurait calculé $V_2=2{,}0\times87/27=\qty{6{,}4}{\litre}$, résultat absurde et faux.
\end{exemple}

\begin{remarque}[variante isochore]
Il existe une troisième situation jumelle, à \textbf{volume constant} (\emph{isochore}) : alors
$p_1/T_1=p_2/T_2$, c'est-à-dire que la pression est proportionnelle à la température absolue. C'est ce
qui fait monter la pression d'un pneu ou d'une bombe aérosol quand ils chauffent. On la déduit
immédiatement de $pV=nRT$ avec $V$ et $n$ constants.
\end{remarque}

% =====================================================================
\subsection{Loi d'Avogadro (volumes égaux \texorpdfstring{$\to$}{->} mêmes nombres de molécules)}

\begin{theoreme}[loi d'Avogadro]
À température $T$ et pression $p$ \textbf{maintenues constantes}, le volume occupé par un gaz est
\textbf{directement proportionnel à sa quantité de matière} :
\[ \boxed{\,\frac{V_1}{n_1}=\frac{V_2}{n_2}\,} \]
Conséquence fondamentale : \emph{des volumes égaux de gaz, pris dans les mêmes conditions de
température et de pression, contiennent le même nombre de molécules}.
\end{theoreme}

\begin{ou}
  \item $V_1,\ n_1$ : volume et quantité de matière du premier gaz (ou état) ;
  \item $V_2,\ n_2$ : volume et quantité de matière du second gaz (ou état) ;
  \item $V$ en unité quelconque (identique), $n$ en moles (\si{\mole}).
\end{ou}

\noindent Cette loi est surprenante : elle dit que la \emph{nature} du gaz n'intervient pas. À $p$ et
$T$ fixés, un litre de \ce{H2} et un litre de \ce{CO2} contiennent \emph{le même nombre de molécules},
bien que ces molécules aient des masses très différentes.

\begin{center}
\begin{tikzpicture}[scale=1.0]
  % récipient 1 : H2
  \draw[thick] (0,0) rectangle (2.6,3);
  \fill[gazcol!7] (0,0) rectangle (2.6,3);
  \foreach \x/\y in {0.5/0.5,1.1/1.2,1.9/0.8,0.8/2.1,2.0/1.9,1.4/2.6,2.2/2.4,0.4/1.6}{
     \fill[azote] (\x,\y) circle (2.6pt);}
  \node[below] at (1.3,-0.25) {$\qty{1}{\litre}$ de \ce{H2}};
  \node[font=\scriptsize] at (1.3,3.35) {8 molécules};
  % égalité conditions
  \node at (3.55,1.5) {\Large $=$};
  \node[font=\scriptsize,align=center] at (3.55,0.7) {même $p$\\ même $T$};
  % récipient 2 : CO2
  \draw[thick] (4.5,0) rectangle (7.1,3);
  \fill[gazcol!7] (4.5,0) rectangle (7.1,3);
  \foreach \x/\y in {5.0/0.5,5.6/1.2,6.4/0.8,5.3/2.1,6.5/1.9,5.9/2.6,6.7/2.4,4.9/1.6}{
     \fill[oxygene] (\x,\y) circle (3.6pt);}
  \node[below] at (5.8,-0.25) {$\qty{1}{\litre}$ de \ce{CO2}};
  \node[font=\scriptsize] at (5.8,3.35) {8 molécules};
  % conclusion
  \node[primarydark,font=\footnotesize,align=center] at (9.6,1.5)
     {\textbf{même volume}\\ $\Rightarrow$ \textbf{même nombre}\\ \textbf{de molécules}\\[2pt]
      (mais masses\\ différentes !)};
\end{tikzpicture}
\end{center}

\begin{exemple}[gonfler un ballon double sa quantité de gaz]
Un ballon contient $n_1=\qty{0{,}50}{\mole}$ de gaz et occupe $V_1=\qty{12{,}2}{\litre}$ à une
certaine température et pression. On y injecte du gaz, toujours aux mêmes $p$ et $T$, jusqu'à atteindre
$n_2=\qty{0{,}80}{\mole}$. Quel volume occupe-t-il ?
\[ \frac{V_1}{n_1}=\frac{V_2}{n_2}
   \;\Longrightarrow\;
   V_2 = V_1\times\frac{n_2}{n_1}
       = \qty{12{,}2}{\litre}\times\frac{0{,}80}{0{,}50}
       = \qty{19{,}52}{\litre}. \]
On a multiplié la quantité de matière par $1{,}6$, le volume est multiplié par le même facteur. Ce
résultat prépare la notion capitale de \textbf{volume molaire} : le volume « par mole » est le même
pour tous les \GP\ dans des conditions données.
\end{exemple}

% =====================================================================
\section{La loi des gaz parfaits : l'équation d'état \texorpdfstring{$pV=nRT$}{pV=nRT}}
% =====================================================================
Les trois lois précédentes ne sont que des cas particuliers. En les fusionnant, on obtient une
\textbf{équation unique} qui relie d'un seul coup les quatre variables d'état : c'est l'\textbf{équation
d'état du \gp}, pierre angulaire de toute la fiche.

\begin{formule}[loi des gaz parfaits]
\[ \boxed{\,p\times V = n\times R\times T\,} \]
\end{formule}

\begin{ou}
  \item $p$ : la pression du gaz, en pascals (\si{\pascal}) si l'on travaille en SI, ou en \si{\atm} /
        \si{\kilo\pascal} selon la valeur de $R$ choisie ;
  \item $V$ : le volume occupé, en mètres cubes (\si{\meter\cubed}) en SI, ou en litres (\si{\litre}) ;
  \item $n$ : la quantité de matière de gaz, en moles (\si{\mole}) ;
  \item $R$ : la \textbf{constante des \GP} (voir ci-dessous) ;
  \item $T$ : la température absolue, \emph{toujours} en kelvins (\si{\kelvin}).
\end{ou}

\begin{remarque}[comment $pV=nRT$ contient les trois lois]
L'équation d'état \emph{résume} tout ce qui précède. Il suffit d'y fixer deux variables :
\begin{itemize}[leftmargin=1.6em,topsep=2pt,itemsep=1.5pt]
  \item à $n$ et $T$ fixés : $pV = (nRT)=\text{cste}\ \Rightarrow$ \textbf{Boyle-Mariotte} ;
  \item à $n$ et $p$ fixés : $V/T = (nR/p)=\text{cste}\ \Rightarrow$ \textbf{Charles--Gay-Lussac} ;
  \item à $p$ et $T$ fixés : $V/n = (RT/p)=\text{cste}\ \Rightarrow$ \textbf{Avogadro}.
\end{itemize}
On n'a donc \emph{qu'une seule} formule à mémoriser : les trois lois historiques en découlent.
\end{remarque}

\subsection{La constante des gaz parfaits \texorpdfstring{$R$}{R}}

\begin{definition}[constante des gaz parfaits]
$R$ est une \textbf{constante universelle} (la même pour tous les gaz). Sa \emph{valeur numérique}
dépend des unités choisies pour $p$ et $V$ ; il faut donc choisir la version de $R$ \textbf{cohérente}
avec ses unités.
\end{definition}

\begin{center}\renewcommand{\arraystretch}{1.5}
\begin{tabular}{c c c c}
\toprule
\rowcolor{primary!12}
\textbf{Valeur de $R$} & \textbf{$p$ en} & \textbf{$V$ en} & \textbf{$T$ en}\\
\midrule
$\qty{8{,}314}{\joule\per\mole\per\kelvin}$           & \si{\pascal}      & \si{\meter\cubed} & \si{\kelvin}\\
$\qty{8{,}314}{\kilo\pascal\litre\per\mole\per\kelvin}$ & \si{\kilo\pascal} & \si{\litre}       & \si{\kelvin}\\
$\qty{0{,}082}{\atm\litre\per\mole\per\kelvin}$        & \si{\atm}         & \si{\litre}       & \si{\kelvin}\\
\bottomrule
\end{tabular}
\end{center}

\begin{attention}
Les deux premières lignes ont la \emph{même} valeur numérique $8{,}314$ car
$\qty{1}{\kilo\pascal}\times\qty{1}{\litre}=\qty{1}{\joule}$. La troisième, $R=\num{0,082}$, est
réservée au couple (\si{\atm}, \si{\litre}) — c'est la plus pratique pour les calculs « chimie » de
cette fiche. \textbf{Ne mélangez jamais} une pression en \si{\atm} avec $R=8{,}314$ : choisissez une
ligne du tableau et tenez-vous-y.
\end{attention}

\subsection{Forme « quantité de matière » et autres réarrangements}

L'équation $pV=nRT$ s'isole selon l'inconnue recherchée. La forme la plus utilisée en chimie donne
directement le \textbf{nombre de moles} à partir des conditions $(p,V,T)$ :

\begin{formule}[nombre de moles d'un gaz]
\[ \boxed{\,n = \dfrac{p\,V}{R\,T}\,} \]
\end{formule}

\begin{ou}
  \item $n$ : la quantité de matière, en moles (\si{\mole}) ;
  \item $p,\ V,\ T$ : pression, volume et température absolue, dans des unités cohérentes avec $R$.
\end{ou}

\noindent On obtient de même, selon le besoin :
\[ p=\frac{nRT}{V},\qquad V=\frac{nRT}{p},\qquad T=\frac{pV}{nR}. \]

\subsection{Conditions normales (CNTP) et volume molaire \texorpdfstring{$\qty{22{,}4}{\litre\per\mole}$}{22,4 L/mol}}

Pour comparer des gaz, on définit des conditions de référence.

\begin{definition}[conditions normales de température et de pression]
Les \textbf{conditions normales de température et de pression} (CNTP, aussi notées TPN ou PTN dans le
livre) correspondent à :
\[ p=\qty{1}{\atm}\qquad\text{et}\qquad T=\qty{273}{\kelvin}\ (\theta=\qty{0}{\degreeCelsius}). \]
\end{definition}

\begin{theoreme}[volume molaire aux CNTP]
Aux CNTP, \textbf{une mole de n'importe quel \gp\ occupe le même volume}, appelé \textbf{volume
molaire} :
\[ \boxed{\,\Vm = \frac{RT}{p} = \qty{22{,}4}{\litre\per\mole}\,} \]
\end{theoreme}

\noindent \textbf{Démonstration.} On part de $pV=nRT$ avec $n=\qty{1}{\mole}$, et l'on calcule le
volume aux CNTP avec $R=\num{0,082}$ :
\[ \Vm = \frac{nRT}{p}
       = \frac{\qty{1}{\mole}\times\num{0,082}\times\qty{273}{\kelvin}}{\qty{1}{\atm}}
       = \qty{22{,}4}{\litre}. \]

\begin{formule}[passage volume $\leftrightarrow$ quantité de matière aux CNTP]
\[ \boxed{\,n = \dfrac{V}{\Vm} = \dfrac{V}{\qty{22{,}4}{\litre\per\mole}}\,}
   \qquad\Longleftrightarrow\qquad
   V = n\times\qty{22{,}4}{\litre\per\mole}. \]
\end{formule}

\begin{ou}
  \item $n$ : quantité de matière de gaz, en moles (\si{\mole}) ;
  \item $V$ : volume du gaz mesuré \emph{aux CNTP}, en litres (\si{\litre}) ;
  \item $\Vm=\qty{22{,}4}{\litre\per\mole}$ : le volume molaire aux CNTP.
\end{ou}

\noindent Le diagramme ci-dessous (repris de la synthèse du livre) résume comment ces relations
s'articulent.

\begin{center}
\begin{tikzpicture}[
   box/.style={draw,thick,rounded corners=2pt,align=center,inner sep=4pt,font=\small},
   lab/.style={font=\scriptsize,midway},
   >=Stealth]
  % bloc gaz parfaits
  \node[box,fill=gazcol!10] (gp) at (0,0) {\textbf{Gaz parfaits}};
  % loi générale
  \node[box,fill=formulacol!12] (gen) at (6.4,0) {Loi générale\\ $n=\dfrac{pV}{RT}$\ (\si{\mole})};
  \draw[->,dashed] (gp) -- node[lab,above]{sont régis par} (gen);
  % CNTP
  \node[box,fill=primary!8,align=left] (cntp) at (1.0,2.6)
     {CNTP : $p=\qty{1}{\atm}$, $T=\qty{273}{\kelvin}$\\
      $R=\qty{0{,}082}{\atm\litre\per\mole\per\kelvin}$\\
      $\dfrac{RT}{p}=\qty{22{,}4}{\litre\per\mole}$};
  % V = 22,4
  \node[box,fill=vcol!12] (v) at (6.4,2.6) {$V=\qty{22{,}4}{\litre}$};
  \draw[->,gazcol] (cntp.east) -- node[lab,above,sloped]{$\div\,22{,}4$}  (v.west);
  \draw[<->] (gen) -- node[lab,right,align=left]{si $n=\qty{1}{\mole}$\\ $\times\,22{,}4$} (v);
  % pression partielle
  \node[box,fill=ncol!10,align=center] (dalton) at (3.2,-2.5)
     {Chaque gaz exerce $p_i=\chi_i\times p_{\text{totale}}$\\[1pt]
      et $p_{\text{totale}}=\displaystyle\sum_i p_i=p_A+p_B+p_C+\dots$\ \ \textbf{(loi de Dalton)}};
  \draw[->,dashed] (gp) -- node[lab,left]{dont} (dalton);
\end{tikzpicture}
\end{center}

\begin{exemple}[du volume au nombre de moles : l'azote aux CNTP]
On dispose de $V=\qty{44{,}8}{\litre}$ de diazote \ce{N2} mesurés aux CNTP. Combien de moles
cela représente-t-il ?

\smallskip
\textbf{Méthode 1 — volume molaire.}
\[ n=\frac{V}{\Vm}=\frac{\qty{44{,}8}{\litre}}{\qty{22{,}4}{\litre\per\mole}}=\qty{2{,}0}{\mole}. \]

\textbf{Méthode 2 — équation d'état (pour vérifier).} Avec $p=\qty{1}{\atm}$, $T=\qty{273}{\kelvin}$,
$R=\num{0,082}$ :
\[ n=\frac{pV}{RT}
    =\frac{\qty{1}{\atm}\times\qty{44{,}8}{\litre}}{\num{0,082}\times\qty{273}{\kelvin}}
    =\frac{44{,}8}{22{,}386}\approx\qty{2{,}0}{\mole}. \]
Les deux méthodes concordent : $\qty{44{,}8}{\litre}$ aux CNTP correspondent à $\qty{2{,}0}{\mole}$.
Le volume molaire $\qty{22{,}4}{\litre\per\mole}$ n'est rien d'autre qu'un raccourci de l'équation
d'état dans les conditions normales.
\end{exemple}

\begin{exemple}[calcul d'une pression : un gaz dans une bouteille]
On enferme $n=\qty{0{,}50}{\mole}$ de dioxygène \ce{O2} dans une bouteille rigide de
$V=\qty{5{,}0}{\litre}$ portée à $\theta=\qty{27}{\degreeCelsius}$. Quelle pression règne à
l'intérieur ?

\smallskip
\textbf{Étape 1 — kelvins :} $T=27+273=\qty{300}{\kelvin}$.

\textbf{Étape 2 — isoler $p$ :} $p=\dfrac{nRT}{V}$.

\textbf{Étape 3 — application numérique} (on prend $R=\num{0,082}$, $V$ en \si{\litre}, donc $p$ sortira
en \si{\atm}) :
\[ p=\frac{\qty{0{,}50}{\mole}\times\num{0,082}\times\qty{300}{\kelvin}}{\qty{5{,}0}{\litre}}
    =\frac{12{,}3}{5{,}0}=\qty{2{,}46}{\atm}. \]
La pression vaut environ $\qty{2{,}5}{\atm}$, soit deux fois et demie la pression atmosphérique.
\end{exemple}

% =====================================================================
\section{Masse volumique, masse molaire et densité d'un gaz}
% =====================================================================
L'équation d'état ne contient pas directement la masse. Mais comme $n=m/M$, on peut y faire entrer la
\textbf{masse} $m$ et la \textbf{masse molaire} $M$ du gaz, et accéder ainsi à sa \textbf{masse
volumique} $\rho$. C'est ce qui permet, à l'inverse, d'\emph{identifier} un gaz inconnu en mesurant sa
masse volumique.

\subsection{Masse volumique d'un gaz}

\begin{definition}[masse volumique]
La \textbf{masse volumique} $\rho$ d'un gaz est la masse contenue dans l'unité de volume :
\[ \rho=\frac{m}{V}. \]
\end{definition}

\begin{ou}
  \item $\rho$ : la masse volumique, en \si{\gram\per\litre} (ou \si{\kilo\gram\per\meter\cubed}) ;
  \item $m$ : la masse de gaz, en grammes (\si{\gram}) ;
  \item $V$ : le volume occupé, en litres (\si{\litre}).
\end{ou}

\noindent \textbf{Établissement de la formule.} On part de $pV=nRT$ et de $n=\dfrac{m}{M}$. En
remplaçant $n$ :
\[ pV=\frac{m}{M}RT
   \;\Longrightarrow\;
   pM=\frac{m}{V}RT=\rho RT. \]
On en tire la masse volumique d'un \gp\ :

\begin{formule}[masse volumique d'un gaz parfait]
\[ \boxed{\,\rho = \dfrac{p\,M}{R\,T}\,} \]
\end{formule}

\begin{ou}
  \item $\rho$ : la masse volumique du gaz, en \si{\gram\per\litre} ;
  \item $p$ : la pression, dans l'unité cohérente avec $R$ ;
  \item $M$ : la masse molaire du gaz, en \si{\gram\per\mole} ;
  \item $R$ : la constante des \GP\ ; \quad $T$ : la température absolue, en \si{\kelvin}.
\end{ou}

\begin{remarque}
Cette formule montre qu'à $p$ et $T$ donnés, la masse volumique d'un gaz est \textbf{proportionnelle à
sa masse molaire} $M$ : un gaz « lourd » (grand $M$, comme \ce{CO2}) est plus dense qu'un gaz « léger »
(petit $M$, comme \ce{He}). C'est pourquoi le \ce{CO2} s'accumule au fond d'une cave et l'hélium fait
monter les ballons.
\end{remarque}

\subsection{Masse molaire d'un gaz}

En isolant $M$ dans la relation $pM=\rho RT$ (ou dans $pV=\frac{m}{M}RT$), on obtient une méthode de
\textbf{détermination expérimentale de la masse molaire} d'un gaz : il suffit de peser une masse $m$ de
gaz occupant un volume $V$ connu dans des conditions $(p,T)$ connues.

\begin{formule}[masse molaire d'un gaz]
\[ \boxed{\,M = \dfrac{m\,R\,T}{p\,V}\,} \]
\end{formule}

\begin{ou}
  \item $M$ : la masse molaire cherchée, en \si{\gram\per\mole} ;
  \item $m$ : la masse de gaz pesée, en grammes (\si{\gram}) ;
  \item $p,\ V,\ T$ : conditions de la mesure (unités cohérentes avec $R$, $T$ en \si{\kelvin}).
\end{ou}

\noindent Aux CNTP, cette relation se simplifie énormément : comme une mole occupe
$\qty{22{,}4}{\litre}$, on a directement
\[ M = \rho\times\Vm = \rho\times\qty{22{,}4}{\litre\per\mole}\qquad(\text{aux CNTP}). \]

\begin{exemple}[identifier un gaz inconnu par sa masse volumique]
Aux CNTP, on mesure la masse volumique d'un gaz inconnu : $\rho=\qty{1{,}25}{\gram\per\litre}$. Quelle
est sa masse molaire, et de quel gaz peut-il s'agir ?

\smallskip
\textbf{Calcul de $M$.} Aux CNTP, $M=\rho\times\Vm$ :
\[ M=\qty{1{,}25}{\gram\per\litre}\times\qty{22{,}4}{\litre\per\mole}=\qty{28}{\gram\per\mole}. \]

\textbf{Identification.} Une masse molaire de $\qty{28}{\gram\per\mole}$ correspond au diazote \ce{N2}
($2\times14=28$) ou au monoxyde de carbone \ce{CO} ($12+16=28$). La mesure de la masse volumique seule
ne permet pas de les distinguer, mais elle restreint déjà fortement les possibilités.

\textbf{Vérification par l'équation complète.} Pour $V=\qty{1}{\litre}$, on a
$m=\rho V=\qty{1{,}25}{\gram}$, et avec $p=\qty{1}{\atm}$, $T=\qty{273}{\kelvin}$ :
\[ M=\frac{mRT}{pV}
    =\frac{\qty{1{,}25}{\gram}\times\num{0,082}\times\qty{273}{\kelvin}}{\qty{1}{\atm}\times\qty{1}{\litre}}
    \approx\qty{28}{\gram\per\mole}. \]
Cohérent.
\end{exemple}

\begin{exemple}[masse volumique du dioxyde de carbone]
Quelle est la masse volumique du \ce{CO2} aux CNTP ? Sa masse molaire est
$M=12+2\times16=\qty{44}{\gram\per\mole}$.
\[ \rho=\frac{M}{\Vm}=\frac{\qty{44}{\gram\per\mole}}{\qty{22{,}4}{\litre\per\mole}}
       =\qty{1{,}96}{\gram\per\litre}. \]
Le \ce{CO2} ($\qty{1{,}96}{\gram\per\litre}$) est nettement plus dense que l'air
($\approx\qty{1{,}29}{\gram\per\litre}$, voir ci-dessous) : il « coule » dans l'air, ce qui explique
son emploi dans les extincteurs (il étouffe la flamme en chassant l'air).
\end{exemple}

\subsection{Composition de l'air et densité d'un gaz par rapport à l'air}

\begin{definition}[composition de l'air]
L'air est un mélange gazeux dont la composition \textbf{en volume} est approximativement :
\[ \approx 80\,\%\ \text{de diazote } \ce{N2} \qquad\text{et}\qquad \approx 20\,\%\ \text{de
dioxygène } \ce{O2}. \]
Sa masse molaire moyenne vaut $M_{\text{air}}=\qty{28{,}8}{\gram\per\mole}$, et $\qty{1}{\litre}$ d'air
a une masse de $\qty{1{,}29}{\gram}$ aux CNTP.
\end{definition}

\begin{center}
\begin{tikzpicture}[scale=1]
  % barre 100% : 80% azote + 20% oxygène
  \fill[azote!75] (0,0) rectangle (8,1);
  \fill[oxygene!75] (8,0) rectangle (10,1);
  \draw[thick] (0,0) rectangle (10,1);
  \node[white] at (4,0.5) {\textbf{\ce{N2} — \SI{80}{\percent}}};
  \node[white] at (9,0.5) {\textbf{\ce{O2}}};
  \node[oxygene,below,font=\scriptsize] at (9,0) {\SI{20}{\percent}};
  \node[below,font=\scriptsize] at (4,-0.05) {composition de l'air \emph{en volume}};
  \node[right,font=\footnotesize,align=left] at (10.4,0.5)
     {$M_{\text{air}}=\qty{28{,}8}{\gram\per\mole}$\\ $\rho_{\text{air}}\approx\qty{1{,}29}{\gram\per\litre}$};
\end{tikzpicture}
\end{center}

\begin{remarque}[d'où vient $\qty{28{,}8}{\gram\per\mole}$ ?]
C'est la moyenne pondérée des masses molaires, selon les proportions volumiques (qui sont aussi des
proportions molaires d'après Avogadro) :
\[ M_{\text{air}} = 0{,}80\times M(\ce{N2}) + 0{,}20\times M(\ce{O2})
   = 0{,}80\times28 + 0{,}20\times32 = 22{,}4+6{,}4 = \qty{28{,}8}{\gram\per\mole}. \]
\end{remarque}

\begin{definition}[densité d'un gaz par rapport à l'air]
La \textbf{densité} $d$ d'un gaz par rapport à l'air est le rapport (sans unité) de sa masse volumique
à celle de l'air, dans les mêmes conditions. Comme à $p,T$ fixés $\rho\propto M$, cela revient au
rapport des masses molaires :
\[ d=\frac{\rho_{\text{gaz}}}{\rho_{\text{air}}}=\frac{M_{\text{gaz}}}{M_{\text{air}}}
    =\frac{M_{\text{gaz}}}{\qty{28{,}8}{\gram\per\mole}}. \]
\end{definition}

\noindent On en déduit la relation du livre, qui donne la masse molaire d'un gaz connaissant sa densité :

\begin{formule}[masse molaire à partir de la densité]
\[ \boxed{\,M_{\text{gaz}} = M_{\text{air}}\times d_{\text{gaz}} = \qty{28{,}8}{\gram\per\mole}\times
   d_{\text{gaz}}\,}\qquad(\text{aux CNTP}). \]
\end{formule}

\begin{ou}
  \item $M_{\text{gaz}}$ : masse molaire du gaz, en \si{\gram\per\mole} ;
  \item $M_{\text{air}}=\qty{28{,}8}{\gram\per\mole}$ : masse molaire moyenne de l'air ;
  \item $d_{\text{gaz}}$ : densité du gaz par rapport à l'air (nombre sans unité).
\end{ou}

\begin{exemple}[ballons : pourquoi l'hélium monte et le \ce{CO2} descend]
Comparons trois gaz à l'air ($M_{\text{air}}=\qty{28{,}8}{\gram\per\mole}$) :
\begin{itemize}[leftmargin=1.6em,topsep=2pt,itemsep=2pt]
  \item \textbf{Hélium} \ce{He}, $M=\qty{4}{\gram\per\mole}$ :
        $d=\dfrac{4}{28{,}8}=0{,}14<1$ $\Rightarrow$ plus léger que l'air, \emph{il monte} (ballons
        gonflés à l'hélium).
  \item \textbf{Dioxyde de carbone} \ce{CO2}, $M=\qty{44}{\gram\per\mole}$ :
        $d=\dfrac{44}{28{,}8}=1{,}53>1$ $\Rightarrow$ plus lourd que l'air, \emph{il descend} (s'accumule
        au sol).
  \item \textbf{Diazote} \ce{N2}, $M=\qty{28}{\gram\per\mole}$ :
        $d=\dfrac{28}{28{,}8}=0{,}97\approx1$ $\Rightarrow$ quasiment la densité de l'air (ce qui est
        logique, l'air étant à $80\,\%$ du \ce{N2}).
\end{itemize}
\textbf{Lecture inverse :} si l'on mesure expérimentalement $d=1{,}53$ pour un gaz inconnu, alors
$M=28{,}8\times1{,}53\approx\qty{44}{\gram\per\mole}$, ce qui oriente vers le \ce{CO2}.
\end{exemple}

% =====================================================================
\section{Mélanges gazeux : la loi de Dalton des pressions partielles}
% =====================================================================
Jusqu'ici, le récipient ne contenait qu'un seul gaz. Or l'air, les gaz de combustion, l'atmosphère
d'un réacteur sont des \textbf{mélanges}. La loi de Dalton décrit comment chaque constituant participe
à la pression totale.

\subsection{Pression partielle}

\begin{definition}[pression partielle]
Dans un mélange gazeux, la \textbf{pression partielle} $p_i$ d'un constituant est la pression qu'il
exercerait \emph{s'il occupait seul} tout le volume du récipient, à la même température. Chaque gaz se
comporte comme si les autres n'existaient pas (hypothèse des \GP\ : pas d'interaction entre molécules).
\end{definition}

\begin{theoreme}[loi de Dalton]
La \textbf{pression totale} d'un mélange de \GP\ est égale à la \textbf{somme des pressions partielles}
de ses constituants :
\[ \boxed{\,p_{\text{tot}} = p_A + p_B + p_C + \dots = \sum_i p_i\,} \]
\end{theoreme}

\begin{ou}
  \item $p_{\text{tot}}$ : la pression totale du mélange, en \si{\atm} (ou autre unité de pression) ;
  \item $p_A,\ p_B,\ p_C,\dots$ : les pressions partielles de chaque gaz, dans la même unité.
\end{ou}

\begin{center}
\begin{tikzpicture}[scale=0.95]
  % récipient A
  \draw[thick] (0,0) rectangle (2.4,2.8);
  \foreach \x/\y in {0.5/0.5,1.1/1.4,1.8/0.8,0.7/2.2,1.9/2.0,1.3/1.0}{
     \fill[azote] (\x,\y) circle (2.8pt);}
  \node[below,font=\footnotesize] at (1.2,-0.2) {\ce{N2} seul};
  \node[azote,above,font=\scriptsize] at (1.2,2.8) {$p_A$};
  % plus
  \node at (3.0,1.4) {\Large $+$};
  % récipient B
  \draw[thick] (3.6,0) rectangle (6.0,2.8);
  \foreach \x/\y in {4.1/0.7,5.0/1.5,5.6/0.6,4.4/2.1,5.3/2.3}{
     \fill[oxygene] (\x,\y) circle (3.4pt);}
  \node[below,font=\footnotesize] at (4.8,-0.2) {\ce{O2} seul};
  \node[oxygene,above,font=\scriptsize] at (4.8,2.8) {$p_B$};
  % egale
  \node at (6.6,1.4) {\Large $=$};
  % récipient mélange (même volume)
  \draw[thick] (7.2,0) rectangle (9.6,2.8);
  \foreach \x/\y in {7.6/0.5,8.2/1.4,8.9/0.8,7.8/2.2,9.0/2.0,8.4/1.0}{
     \fill[azote] (\x,\y) circle (2.8pt);}
  \foreach \x/\y in {8.6/0.5,7.7/1.5,9.2/1.4,8.1/2.3,8.9/2.4}{
     \fill[oxygene] (\x,\y) circle (3.4pt);}
  \node[below,font=\footnotesize] at (8.4,-0.2) {mélange};
  \node[primarydark,above,font=\scriptsize] at (8.4,2.8) {$p_{\text{tot}}=p_A+p_B$};
\end{tikzpicture}
\end{center}

\noindent\textbf{Origine de la loi.} Chaque gaz suit sa propre équation d'état dans le volume commun
$V$ à la température commune $T$ : $p_i V = n_i RT$. En sommant sur tous les gaz :
\[ \Big(\sum_i p_i\Big)V = \Big(\sum_i n_i\Big)RT = n_{\text{tot}}RT, \]
ce qui montre que $\sum_i p_i$ joue le rôle de la pression totale : $p_{\text{tot}}V=n_{\text{tot}}RT$.

\subsection{Fraction molaire et pression partielle}

\begin{definition}[fraction molaire]
La \textbf{fraction molaire} $\chi_A$ d'un constituant $A$ est la proportion de ses moles parmi le
nombre total de moles du mélange :
\[ \chi_A = \frac{n_A}{n_{\text{total}}}. \]
C'est un nombre \emph{sans unité}, compris entre $0$ et $1$. La somme des fractions molaires de tous
les constituants vaut $1$ : $\sum_i \chi_i = 1$.
\end{definition}

\begin{ou}
  \item $\chi_A$ : la fraction molaire du gaz $A$ (sans unité) ;
  \item $n_A$ : la quantité de matière du gaz $A$, en moles (\si{\mole}) ;
  \item $n_{\text{total}}$ : la quantité de matière totale du mélange, en moles (\si{\mole}).
\end{ou}

\begin{formule}[pression partielle et fraction molaire]
\[ \boxed{\,p_A = \chi_A \times p_{\text{totale}}\,} \]
\end{formule}

\begin{ou}
  \item $p_A$ : la pression partielle du gaz $A$, en \si{\atm} ;
  \item $\chi_A$ : la fraction molaire du gaz $A$ (sans unité) ;
  \item $p_{\text{totale}}$ : la pression totale du mélange, en \si{\atm}.
\end{ou}

\noindent\textbf{Justification.} En divisant $p_A V = n_A RT$ par $p_{\text{tot}}V=n_{\text{tot}}RT$,
les facteurs $V$, $R$, $T$ se simplifient :
\[ \frac{p_A}{p_{\text{tot}}}=\frac{n_A}{n_{\text{tot}}}=\chi_A
   \quad\Longrightarrow\quad p_A=\chi_A\,p_{\text{tot}}. \]
La pression partielle d'un gaz est donc, tout simplement, sa \emph{part} de la pression totale —
proportionnelle à sa fraction molaire.

\begin{remarque}
Pour des \GP, la fraction \textbf{molaire} est égale à la fraction \textbf{volumique} (conséquence de
la loi d'Avogadro). C'est pourquoi « l'air contient $20\,\%$ de \ce{O2} en volume » signifie aussi
$\chi(\ce{O2})=0{,}20$, et donc $p(\ce{O2})=0{,}20\times p_{\text{atm}}$.
\end{remarque}

\begin{exemple}[pressions partielles dans un mélange \ce{N2}/\ce{O2}]
Un récipient rigide de $V=\qty{10{,}0}{\litre}$ à $T=\qty{300}{\kelvin}$ contient
$n(\ce{N2})=\qty{0{,}30}{\mole}$ et $n(\ce{O2})=\qty{0{,}20}{\mole}$. Déterminons la pression totale et
les pressions partielles.

\smallskip
\textbf{Étape 1 — quantité totale.}
$n_{\text{tot}}=0{,}30+0{,}20=\qty{0{,}50}{\mole}$.

\textbf{Étape 2 — pression totale} (équation d'état, $R=\num{0,082}$) :
\[ p_{\text{tot}}=\frac{n_{\text{tot}}RT}{V}
   =\frac{\qty{0{,}50}{\mole}\times\num{0,082}\times\qty{300}{\kelvin}}{\qty{10{,}0}{\litre}}
   =\frac{12{,}3}{10{,}0}=\qty{1{,}23}{\atm}. \]

\textbf{Étape 3 — fractions molaires.}
\[ \chi(\ce{N2})=\frac{0{,}30}{0{,}50}=0{,}60,\qquad \chi(\ce{O2})=\frac{0{,}20}{0{,}50}=0{,}40
   \qquad(\text{somme}=1\ \checkmark). \]

\textbf{Étape 4 — pressions partielles.}
\[ p(\ce{N2})=0{,}60\times1{,}23=\qty{0{,}74}{\atm},\qquad
   p(\ce{O2})=0{,}40\times1{,}23=\qty{0{,}49}{\atm}. \]

\textbf{Vérification (loi de Dalton).}
$p(\ce{N2})+p(\ce{O2})=0{,}74+0{,}49=\qty{1{,}23}{\atm}=p_{\text{tot}}$ $\checkmark$. Tout est
cohérent.
\end{exemple}

\begin{exemple}[combien d'oxygène dans une pièce ? une application de l'air]
Combien de moles — et quelle masse — de dioxygène \ce{O2} contient une pièce de dimensions
$\qty{2}{\meter}\times\qty{5{,}6}{\meter}\times\qty{2}{\meter}$ aux CNTP, sachant que l'air contient
$20\,\%$ de \ce{O2} en volume ?

\smallskip
\textbf{Étape 1 — volume d'air.}
\[ V_{\text{air}}=2\times5{,}6\times2=\qty{22{,}4}{\meter\cubed}=\qty{22400}{\litre}. \]

\textbf{Étape 2 — moles d'air total} (aux CNTP, $\Vm=\qty{22{,}4}{\litre\per\mole}$) :
\[ n_{\text{air}}=\frac{V_{\text{air}}}{\Vm}=\frac{\qty{22400}{\litre}}{\qty{22{,}4}{\litre\per\mole}}
   =\qty{1000}{\mole}. \]

\textbf{Étape 3 — moles de \ce{O2}} (fraction volumique $=$ fraction molaire $=0{,}20$) :
\[ n(\ce{O2})=0{,}20\times1000=\qty{200}{\mole}. \]

\textbf{Étape 4 — masse de \ce{O2}} ($M(\ce{O2})=2\times16=\qty{32}{\gram\per\mole}$) :
\[ m(\ce{O2})=n\times M=200\times32=\qty{6400}{\gram}=\qty{6{,}4}{\kilo\gram}. \]
La pièce contient $\qty{200}{\mole}$, soit $\qty{6{,}4}{\kilo\gram}$, de dioxygène. Remarquez l'astuce
de l'énoncé : le volume tombe pile sur $\qty{22{,}4}{\meter\cubed}$, donc $\qty{1000}{\mole}$ d'air —
un clin d'œil au volume molaire.
\end{exemple}

\begin{exemple}[passer d'un volume de gaz à une masse via la stœchiométrie]
\emph{Cet exemple relie la loi des gaz à un bilan de matière (Fiche 5).} On fait réagir
$\qty{44{,}8}{\litre}$ de \ce{N2} et $\qty{67{,}2}{\litre}$ de \ce{H2}, mesurés aux CNTP, selon la
synthèse de l'ammoniac :
\[ \ce{N2(g) + 3 H2(g) -> 2 NH3(g)}. \]
Quelle masse maximale d'ammoniac \ce{NH3} peut-on obtenir ?

\smallskip
\textbf{Étape 1 — volumes $\to$ moles} (aux CNTP, on divise par $\qty{22{,}4}{\litre\per\mole}$) :
\[ n(\ce{N2})=\frac{44{,}8}{22{,}4}=\qty{2{,}0}{\mole},\qquad
   n(\ce{H2})=\frac{67{,}2}{22{,}4}=\qty{3{,}0}{\mole}. \]

\textbf{Étape 2 — réactif limitant.} La réaction consomme $1$ \ce{N2} pour $3$ \ce{H2}. Pour
$\qty{2{,}0}{\mole}$ de \ce{N2}, il faudrait $3\times2=\qty{6{,}0}{\mole}$ de \ce{H2} ; on n'en a que
$\qty{3{,}0}{\mole}$. C'est donc le \textbf{dihydrogène \ce{H2} qui est limitant}.

\textbf{Étape 3 — quantité d'ammoniac formé.} D'après les coefficients, $3$ \ce{H2} donnent $2$ \ce{NH3} :
\[ n(\ce{NH3})=\frac{2}{3}\times n(\ce{H2})=\frac{2}{3}\times3{,}0=\qty{2{,}0}{\mole}. \]

\textbf{Étape 4 — masse} ($M(\ce{NH3})=14+3\times1=\qty{17}{\gram\per\mole}$) :
\[ m(\ce{NH3})=n\times M=2{,}0\times17=\qty{34}{\gram}. \]
On obtient au mieux $\qty{34}{\gram}$ d'ammoniac. Cet exemple montre la \textbf{chaîne de raisonnement}
typique en chimie des gaz : \emph{volume} $\xrightarrow{\,\div\,22{,}4\,}$ \emph{moles}
$\xrightarrow{\text{stœchiométrie}}$ \emph{moles du produit} $\xrightarrow{\,\times M\,}$ \emph{masse}.
\end{exemple}

% =====================================================================
\section{Synthèse}
% =====================================================================

\noindent Le schéma suivant relie l'ensemble des notions de la fiche, du cas particulier (lois
historiques) au cas général (équation d'état), puis aux mélanges (Dalton).

\begin{center}
\begin{tikzpicture}[
   node distance=6mm,
   b/.style={draw,thick,rounded corners=3pt,align=center,inner sep=5pt,font=\small,fill=white},
   bgp/.style={b,fill=gazcol!12},
   bf/.style={b,fill=formulacol!12},
   bd/.style={b,fill=ncol!10},
   >=Stealth]
  \node[bf] (etat) {\textbf{Équation d'état du gaz parfait}\\[2pt] $pV=nRT$};
  \node[bgp,below left=12mm and 6mm of etat] (boyle) {\textbf{Boyle-Mariotte}\\ ($T$ cste)\\ $p_1V_1=p_2V_2$};
  \node[bgp,below=12mm of etat] (charles) {\textbf{Charles--Gay-Lussac}\\ ($p$ cste)\\ $\dfrac{V_1}{T_1}=\dfrac{V_2}{T_2}$};
  \node[bgp,below right=12mm and 6mm of etat] (avo) {\textbf{Avogadro}\\ ($p,T$ cstes)\\ $\dfrac{V_1}{n_1}=\dfrac{V_2}{n_2}$};
  \draw[->] (etat) -- (boyle);
  \draw[->] (etat) -- (charles);
  \draw[->] (etat) -- (avo);
  % branche masse / densité
  \node[bf,right=14mm of etat] (masse) {\textbf{Masse / densité}\\ $\rho=\dfrac{pM}{RT}$,\quad $M=\dfrac{mRT}{pV}$\\ $M_{\text{gaz}}=28{,}8\,d$};
  \draw[->] (etat) -- (masse);
  % branche CNTP
  \node[bf,above=10mm of etat] (cntp) {\textbf{CNTP} : $\qty{1}{\atm}$, $\qty{273}{\kelvin}$\\ $\Vm=\qty{22{,}4}{\litre\per\mole}$,\quad $n=\dfrac{V}{22{,}4}$};
  \draw[->] (etat) -- (cntp);
  % branche Dalton
  \node[bd,below=12mm of charles] (dalton) {\textbf{Mélanges — loi de Dalton}\\ $p_{\text{tot}}=\sum p_i$,\quad $\chi_A=\dfrac{n_A}{n_{\text{tot}}}$,\quad $p_A=\chi_A\,p_{\text{tot}}$};
  \draw[->] (charles) -- (dalton);
\end{tikzpicture}
\end{center}

\begin{tcolorbox}[boxbase,colback=formulacol!6,colframe=formulacol,
   title={\textsc{Mémo — les formules à connaître}}]
\begin{center}\renewcommand{\arraystretch}{1.5}
\begin{tabular}{l l l}
\toprule
\textbf{Situation} & \textbf{Formule} & \textbf{Conditions}\\
\midrule
Cas général                 & $pV=nRT$                       & toujours (gaz parfait)\\
$T$ constante               & $p_1V_1=p_2V_2$                & Boyle-Mariotte\\
$p$ constante               & $V_1/T_1=V_2/T_2$              & Charles--Gay-Lussac\\
$p,T$ constantes            & $V_1/n_1=V_2/n_2$             & Avogadro\\
Nombre de moles             & $n=pV/RT$                      & toujours\\
Aux CNTP                    & $n=V/22{,}4$                  & $p=\qty{1}{\atm}$, $T=\qty{273}{\kelvin}$\\
Masse volumique             & $\rho=pM/(RT)$                & toujours\\
Masse molaire               & $M=mRT/(pV)$                 & toujours\\
Densité / air               & $M_{\text{gaz}}=28{,}8\,d$    & aux CNTP\\
Pression partielle          & $p_A=\chi_A\,p_{\text{tot}}$  & mélange (Dalton)\\
Pression totale             & $p_{\text{tot}}=\sum p_i$     & mélange (Dalton)\\
\bottomrule
\end{tabular}
\end{center}
\end{tcolorbox}

\begin{tcolorbox}[boxbase,colback=primary!6,colframe=primary,
   title={\textsc{Mémo — valeurs et constantes de référence}}]
\begin{itemize}[leftmargin=1.6em,topsep=2pt,itemsep=2pt]
  \item \textbf{Constante des gaz} : $R=\qty{8{,}314}{\joule\per\mole\per\kelvin}=\qty{8{,}314}{\kilo\pascal\litre\per\mole\per\kelvin}=\qty{0{,}082}{\atm\litre\per\mole\per\kelvin}$.
  \item \textbf{CNTP} (= TPN = PTN) : $p=\qty{1}{\atm}$, $T=\qty{273}{\kelvin}$ ($\theta=\qty{0}{\degreeCelsius}$).
  \item \textbf{Volume molaire aux CNTP} : $\Vm=\qty{22{,}4}{\litre\per\mole}$.
  \item \textbf{Air} : $\approx 80\,\%$ \ce{N2} $+$ $20\,\%$ \ce{O2} (en volume) ; $M_{\text{air}}=\qty{28{,}8}{\gram\per\mole}$ ; $\rho_{\text{air}}\approx\qty{1{,}29}{\gram\per\litre}$ (CNTP).
  \item \textbf{Conversions} : $\qty{1}{\atm}=\qty{101325}{\pascal}=\qty{760}{\mmHg}$ ; $\qty{1}{\meter\cubed}=\qty{1000}{\litre}$ ; $T(\si{\kelvin})=\theta+273{,}15$.
\end{itemize}
\end{tcolorbox}

\begin{tcolorbox}[boxbase,colback=attncol!6,colframe=attncol,
   title={\textsc{Mémo — les pièges classiques à éviter}}]
\begin{itemize}[leftmargin=1.6em,topsep=2pt,itemsep=1.5pt]
  \item \textbf{Température} : toujours en \textbf{kelvins} dans $pV=nRT$ et dans les rapports $V/T$, $p/T$. Convertir avant tout calcul.
  \item \textbf{Cohérence $R$ / unités} : associer $R=\num{0,082}$ à (\si{\atm}, \si{\litre}), ou $R=8{,}314$ à (\si{\pascal}, \si{\meter\cubed}). Ne jamais panacher.
  \item \textbf{Le volume $\qty{22{,}4}{\litre\per\mole}$} n'est valable qu'\textbf{aux CNTP}. Hors CNTP, repasser par $pV=nRT$.
  \item \textbf{Fraction molaire $=$ fraction volumique} pour les \GP\ seulement (grâce à Avogadro).
  \item \textbf{Pression partielle} $p_A=\chi_A\,p_{\text{tot}}$ : c'est la fraction \emph{molaire}, pas la fraction massique.
\end{itemize}
\end{tcolorbox}

\begin{tcolorbox}[boxbase,colback=methcol!8,colframe=methcol,sharp corners,
   title={\textsc{Méthode générale — résoudre un problème de gaz}}]
\begin{enumerate}[leftmargin=1.8em,topsep=2pt,itemsep=2pt]
  \item \textbf{Lister} les grandeurs connues et l'inconnue ; repérer ce qui est \emph{constant}.
  \item \textbf{Convertir} les unités : température en \si{\kelvin}, pression et volume cohérents avec le $R$ choisi.
  \item \textbf{Choisir la loi} : une variable constante $\to$ loi historique ; sinon $\to$ $pV=nRT$ ; mélange $\to$ Dalton.
  \item \textbf{Isoler} littéralement l'inconnue \emph{avant} l'application numérique.
  \item \textbf{Calculer}, puis \textbf{vérifier} l'ordre de grandeur et l'homogénéité (les unités doivent se simplifier correctement).
\end{enumerate}
\end{tcolorbox}

\vfill
\begin{center}
\textcolor{primary}{\rule{0.5\textwidth}{1pt}}\\[4pt]
{\small\color{black!60}Fin de la Fiche 7 — Loi des gaz.}
\end{center}

\end{document}
